\setcounter{section}{5}
\section{Make three recommendations}
\label{sec:recommendations}

\paragraph{Rapid classroom demonstration.}
Use Setting~B Q-learning ($\alpha=.15$, $\gamma=.995$, $\epsilon_{\mathrm{end}}=.05$) with a
prepared replay. Evidence: (1) Q trains in 3.8\,min ($225.5\pm 17.9$\,s) and led 6-epoch PPO at
65,536 steps (5.6\% vs.\ 0\%, \texttt{3142fac3111db314}); (2) table entries directly explain
choices (TS1/166: NOOP 0.79 vs.\ LEFT 0.00, +200; TS1/437: near-zero pinch values). Weakness:
68.7\% failure rate. Switch to frozen PPO12 (90.2\%) if live target completion takes priority.
Before non-simulation classroom use, pilot whether students accurately explain Q-value ties on
unseen replays.

\paragraph{Repeated operation in uncertainty.}
Recommend \emph{neither} workflow for consequential autonomous deployment; for further B
simulation, use Part~3's \textbf{PPO12} ($2\times$ LR, 12 epochs) without C's penalty. Evidence:
(1) PPO12 achieved $90.2\pm 8.0\%$ success (9.8\% failure) vs.\ 70.2\% baseline PPO and 31.3\% Q
(Table~\ref{tab:part3-heldout}); (2) C's penalty cut baseline PPO by 3.6\,pp and widened SD to
26.3\,pp, opposing $F>R>E>T$. Weakness: the weakest PPO12 seed still failed 22.5\% of episodes,
and SF2 failed at 1,990 points. Reverse only when fresh retrainings stay within a prespecified
failure bound. Before deployment, test controlled action latencies, hazard severity and human
overrides on a hardware-in-the-loop testbed.

\paragraph{Operation after an environment change.}
Recommend \emph{neither} workflow for unfamiliar conditions. Evidence: (1) under D's fixed
32$\times$32 obstruction, baseline PPO fell 22.7\,pp ($70.2\%\to 47.5\%$) and PPO12 fell 21.0\,pp
($91.1\%\to 70.1\%$, 29.9\% failure); (2) Q dropped less ($-4.9$\,pp) but still failed 77.8\%
(Table~\ref{tab:robustness-pairs}). Weakness: one small blind spot unpredictably redirects seeds.
Reverse only if locked policies meet required reliability across multiple independent shifts.
Before transfer, test unseen obstruction locations, variable dropout durations and pixel
perception on a protected testbed.

\paragraph{Scope and broader principles.}
Results reflect RAM-derived features, sticky dynamics and a 262,144-step cap, not universal
algorithm superiority. Broader principles illustrated are that directional function approximation
scales better than coarse tabular aliasing, rollout reuse governs PPO's early sample efficiency,
and high source-domain success does not certify observation robustness.
